\newpage

\subsection{Development Cards}

In order to build and test the real time software, dedicated development hardware is required to ensure an isolated testing environment.

These test cards consist of a Raspberry Pi Compute Module 4 (CM4).
This contains the same overall hardware as the standard raspberry pi 4.
EMMC flash for high vibration environments.
Exception being the ethernet IC and EMMC flash (This later caused problems with compiling the operating system and bootloader)

Key requirement is remote debugging, full remote power and data control.
Also placed in a rack
Selected the Pi KVM project for remote control as is supports deploy as well as additional GPIO control.

The CM4s display is captured using the CSI port on the KVM with a HDMI to CSI converter.
The power is controlled using a relay, that controls the 12v supply to the CM4.
The signal flags of the EMMC storage are directly connected to the KVM. EXPLAIN the signal flags
These are configured in the KVM user interface
UART is connected to a USB serial converter.
USB host is connected to the KVM for flashing the CM4.
Finally a debugger (JTAG) in case of modifications to the kernel itself.

\begin{figure}[ht]
    \centering
        \begin{tikzpicture}  
            \node[block, text width=1.75cm ] (h) {HDMI\\Capture};  
            \node[block, text width=1.75cm, right=0.1cm of h] (u) {USB\\Serial};  
            \node[block, text width=1.75cm, right=0.1cm of u] (j) {JTAG\\Debugger};  
            \node[block, text width=1.75cm, right=0.1cm of j] (p) {Power\\Relay};

            \node[block, text width=1.75cm, right=0.1cm of p] (m) {Flash\\Mode};
            \node[block, text width=1.75cm, right=0.1cm of m] (l) {Flash\\Lock};
            \node[block, text width=1.75cm, right=0.1cm of l] (f) {USB};

            \node[block, text width=11.75cm, above=1cm of p] (d) {Development Card};  
            \node[block, text width=11.75cm, below=1cm of p] (k) {Remote KVM};  

            \draw[<->,thick] ([xshift=-4.5cm]d.south) to node [midway,fill=white] {HDMI} (h.north);
            \draw[<->,thick] ([xshift=-3cm]d.south) to node [midway,fill=white] {UART} (u.north);
            \draw[<->,thick] ([xshift=-1.5cm]d.south) to node [midway,fill=white] {JTAG} (j.north);
            \draw[<->,thick] ([xshift=-0cm]d.south) to node [midway,fill=white] {12V} (p.north);
            \draw[<->,thick] ([xshift=1.5cm]d.south) edge (m.north);
            \draw[<->,thick] ([xshift=3cm]d.south) edge (l.north);
            \draw[<->,thick] ([xshift=4.5cm]d.south) edge (f.north);

            \draw[<->,thick] (h.south) to node [midway,fill=white] {CSI} ([xshift=-4.5cm]k.north);
            \draw[<->,thick] (u.south) to node [midway,fill=white] {USB} ([xshift=-3cm]k.north);
            \draw[<->,thick] (j.south) to node [midway,fill=white] {USB} ([xshift=-1.5cm]k.north);
            \draw[<->,thick] (p.south) to node [midway,fill=white] {GPIO} ([xshift=0cm]k.north);
            \draw[<->,thick] (m.south) to node [midway,fill=white] {GPIO} ([xshift=1.5cm]k.north);
            \draw[<->,thick] (l.south) to node [midway,fill=white] {GPIO} ([xshift=3cm]k.north);
            \draw[<->,thick] (f.south) to node [midway,fill=white] {USB} ([xshift=4.5cm]k.north);
        \end{tikzpicture}  
    \caption[Development Hardware]
    {Development Hardware}
    \label{fig:developmentHardware}
\end{figure}


Two of these so that multiple benchmarking activities can be performed at once and compared.


Use the uBoot boot loader.
This allows the development images to be loaded over TFTP, we have a dedicated tftp boot server.

Connected to a managed switch. The test harness is connected to this on a separate VLAN.

% \begin{figure}
%     \centering
%     \resizebox{\columnwidth}{!}{
%         %
%         \begin{tikzpicture}  
%             \node[block] (d1) {Development Card};  
%             \node[block,right=of a] (k1) {Debug KVM};

%             \node[block,below=of a] (d) {Reciever};  
%             \node[block,right=of d] (e) {Destination};  

%             % \node[draw,inner xsep=4mm,inner ysep=7mm,fit=(a)(b) ] (c1) {Application Layer};
%             % \node[draw,inner xsep=4mm,inner ysep=7mm,fit=(d)(e) ] {Infrastructure Layer};  
%             % \node[draw,inner xsep=4mm,inner ysep=7mm,fit=(d)(e),label={Hardware Layer}]{};
%         \end{tikzpicture}  
%         %
%     }
%     \caption[Heatmap for class distribution in dataset]
%     {Heatmap for class distribution in dataset}
%     \label{fig:heatmap}
% \end{figure}

\begin{tikzpicture}[node distance=1.75cm]
    \node[draw] (node1) {Node 1};
    \node[draw, below of=node1] (node2) {Node 2};
    
    % Create a node that fits both node1 and node2
    \node[fit=(node1)(node2), draw, inner sep=0pt] (fitnode) {};
    
    % Adjusting the fitting node to match the sizes
    \node[above right] at (fitnode.south west) {Fitting Node};
\end{tikzpicture}

\begin{tikzpicture}[
    mymatrix/.style={matrix of nodes, nodes=block, row sep=1em},
    mycontainer/.style={draw=gray, inner sep=1ex},
    typetag/.style={draw=gray, inner sep=1ex, anchor=west},
    title/.style={draw=none, color=gray, inner sep=0pt}
    ]
    \matrix[mymatrix] (mx1) {
      |[title]|Decomposition \\
      Independent \\
      Reduction \\
      DivideAndConquer \\
    };
    \matrix[mymatrix, right=of mx1.north east, matrix anchor=north west] (mx2) {
      |[title]|Dependency \\
      Atomic \\
      Range \\
    };
    \node[mycontainer, fit=(mx1)] {};
    \node[mycontainer, fit=(mx2)] {};
  
  \end{tikzpicture}